\section{Discussion}

\subsection{Limitations of Our Operationalization}

We must be transparent: our regex-based extraction achieved only 49.5\% coverage, failing to meet the 95\% target. This limits our ability to draw conclusions about AS predictive power. The primary cause is the dominance of assertion errors (53.4\%), which produce output like \texttt{AssertionError} without file:line information.

\subsection{What We Can Conclude}

Despite the extraction limitation, several findings are informative:

\begin{enumerate}
    \item \textbf{Error type stratification is necessary}: Future AS research must separate error types before measuring.
    \item \textbf{AS ordering holds where measurable}: On the extractable subset, C1$\geq$C2$\geq$C3$\geq$C4 ordering is preserved.
    \item \textbf{The framework identifies the barrier}: AS decomposition reveals \emph{why} extraction fails (missing file:line structure) and suggests solutions.
\end{enumerate}

\subsection{Actionable Guidance from the AS Framework}

The AS framework provides more than taxonomy. It identifies:
\begin{itemize}
    \item \textbf{Which errors block extraction}: Assertion errors lacking traceback frames
    \item \textbf{Why they block extraction}: Missing file:line format that extraction patterns require
    \item \textbf{How to address it}: Use \texttt{pytest --tb=long} for richer tracebacks, or LLM-based extraction that doesn't rely on regex patterns
\end{itemize}

\subsection{Limitations}

\begin{itemize}
    \item \textbf{Single operationalization tested} (regex extraction only)
    \item \textbf{Single benchmark family} (HumanEval/MBPP)
    \item \textbf{Python-specific} trace mechanisms
    \item \textbf{Did not test AS predictive power} (blocked by extraction failure)
\end{itemize}
